\documentclass[11pt,a4paper]{article}

\usepackage[margin=2.2cm]{geometry}
\usepackage{amsmath,amssymb}
\usepackage{booktabs}
\usepackage{caption}
\usepackage{enumitem}

\setlist{nosep,leftmargin=1.3em}
\captionsetup{font=small,labelfont=bf}
\renewcommand{\arraystretch}{1.15}
\setlength{\parskip}{0.45em}
\setlength{\parindent}{0pt}

\title{\vspace{-1.4cm}\textbf{Planned Experiments for the Attention-Map Segmentation Stage}}
\author{Narmeen Sabah Siddiqui, Faizyab Ali Shah}
\date{}

\begin{document}

\maketitle
\vspace{-0.8em}

\section*{Planned segmentation experiments}

In the next phase, we will focus on attention-based semantic segmentation of the
generated urban scenes. Rather than treating this as a single-method comparison,
we plan to study how semantic localization, spatial refinement, attention resolution,
and denoising timestep selection affect mask quality.

The experiments are organized in two groups. We first prioritize experiments that
can mainly be carried out using simpler operations on attention information extracted during image
generation. Cross-attention and self-attention maps can be stored at different
resolutions, layers, heads, and denoising timesteps, allowing several ablations and
combinations to be evaluated afterwards without requiring a substantially different
generation pipeline.

After these experiments, we will move to methods that require additional
optimization, more complex operations, external models, or more specialized segmentation procedures. The
results of the first group will also help determine which configurations are worth
using as inputs to these more complex methods.

% =========================================================
\subsection*{Attention-map experiments and ablations}
% =========================================================

\textbf{OVAM baseline.}
OVAM remains the main cross-attention baseline. The current baseline aggregates
cross-attention information across layers, heads, and timesteps before converting
the resulting heatmap into a semantic mask. This configuration provides the
reference against which the following attention-map experiments will be compared.

For the baseline, we will keep the best post-processing configuration found in
Phase~1 and, where necessary, revisit normalization, thresholding, and background
handling. Once fixed, these settings should remain unchanged when comparing the
subsequent attention strategies.

\textbf{OVAM with self-attention refinement.}
Cross-attention provides class-specific semantic information, whereas
self-attention describes relationships between spatial locations. Instead of using
the OVAM cross-attention heatmap directly as the final mask, we will investigate
whether self-attention can propagate or refine the semantic response spatially.

Because the self-attention maps can be stored during generation, several variants
can be evaluated afterwards. The main ablations will include the self-attention
resolution used for refinement ($16\times16$, $32\times32$, and $64\times64$),
the strength of the refinement, and whether the refinement is applied using one
resolution or progressively across multiple resolutions.

A simple cross-attention/self-attention combination will first be evaluated before
moving to the more specialized refinement mechanisms described later.

\textbf{Resolution-specific cross-attention.}
OVAM normally aggregates attention maps originating from several spatial
resolutions. However, the information contained at each resolution is different.
Lower-resolution maps may provide stronger semantic localization, while
higher-resolution maps may contain finer spatial details.

We will therefore store and evaluate the $16\times16$, $32\times32$, and
$64\times64$ cross-attention maps separately, followed by combinations such as
$16+32$, $16+64$, $32+64$, and all resolutions together.

The main objective is not only to determine the best resolution overall, but also
to determine whether this varies between classes. For example, large classes such
as road, sky, building, or vegetation may benefit from coarse semantic attention,
whereas person, pole, traffic sign, or traffic light may require finer-resolution
information.

\textbf{Denoising timestep selection.}
The current OVAM formulation aggregates attention information across the diffusion
trajectory. However, semantic and spatial information may not be equally useful at
all denoising stages.

Since the attention maps can be stored independently for each diffusion timestep,
we will compare the complete trajectory against selected timestep ranges. Initial
experiments will consider early, middle, and late windows, followed by smaller
windows or individual timesteps if clear trends appear.

For example, the experiment may compare:
\[
\text{all timesteps},
\qquad
\text{early},
\qquad
\text{middle},
\qquad
\text{late}.
\]

If necessary, this can later be extended to combinations of resolution and timestep,
so that the attention used for a class can be written conceptually as
\[
A^{c}_{r,t},
\]
where $c$ denotes the class, $r$ the spatial resolution, and $t$ the selected
timestep or timestep window.

\textbf{Class- or scale-adaptive attention strategy.}
The previous experiments can naturally lead to a simple adaptive strategy. Instead
of selecting one global OVAM configuration for every class, each class
can use the attention configuration that performs best for its spatial characteristics.

For example, if the resolution and timestep experiments show that road performs
best with coarse attention while traffic signs perform better with higher-resolution
maps, those configurations can be selected independently.

An initial adaptive strategy will therefore compare:
\[
\text{one global attention configuration}
\]
against
\[
\text{class-specific or scale-specific attention configurations}.
\]

This first version can be implemented entirely using the stored attention maps.
More complex refinement methods such as SAM can later be incorporated if they
prove useful.

\vspace{0.5em}
\hrule
\vspace{0.8em}

% =========================================================
\subsection*{More method-specific extensions}
% =========================================================

The following experiments require additional optimization, an external segmentation
model, or a more specialized processing pipeline. They will therefore be considered
after the basic attention-map experiments above have established which attention
sources are most useful for Cityscapes.

\textbf{Above ablations with token optimization.}
The current OVAM experiments query classes using their ordinary text embeddings,
for example ``car'', ``road'', or ``traffic sign''. OVAM also provides a token
optimization procedure in which the embedding representing a class is optimized
using a very small amount of annotated data, while keeping the diffusion model
itself frozen.

We will compare ordinary class-name embeddings against optimized class tokens and
measure the improvement separately for eachclass. This experiment can
also help distinguish semantic failures from spatial-resolution failures. If a class
remains poorly localized after token optimization, the limitation is more likely to
come from the spatial attention representation rather than from the text embedding.

\textbf{Above ablations with SAM-based refinement.}
In this strategy, diffusion cross-attention is used primarily for semantic localization,
while a segmentation foundation model such as SAM is responsible for obtaining the
final spatial boundary.

The response can be converted into different types of SAM prompts. We will
compare high-confidence points, multiple points, bounding boxes, and coarse mask
prompts. The purpose is to determine how much information should be transferred
from diffusion attention to SAM.

Particular attention will be given to smaller classes.

\textbf{SeeDiff (2025).}
SeeDiff uses a seeded-segmentation formulation. Instead of treating the
cross-attention response as the final segmentation mask, it extracts
high-confidence cross-attention locations as semantic seeds.

These seeds are subsequently expanded using self-attention affinities. The
expansion proceeds progressively from coarser self-attention maps towards
higher-resolution maps, combining semantic consistency at coarse resolutions with
finer spatial information at higher resolutions. A separate background expansion
stage is then used to suppress incorrectly expanded foreground regions.

The main factors to investigate will include the cross-attention seed threshold,
the self-attention resolutions used during expansion, whether expansion is
single-stage or progressive, and the effect of the background refinement.

\textbf{Hierarchical or weighted attention fusion.}
The previous ablations treat attention maps explicitly and allow us to identify
which resolutions, timesteps, and potentially layers are most useful. A more
advanced alternative is to combine this information through weighted or
hierarchical fusion rather than a simple average.

FA-Seg motivates a hierarchical multi-resolution combination of cross-attention
and self-attention, where different scales contribute differently to semantic
localization and spatial refinement. NERVE provides a different motivation by
using uncertainty or entropy to avoid treating every self-attention map equally.

We will therefore consider comparisons such as
\[
A_{\mathrm{uniform}}
=
\frac{1}{N}\sum_i A_i
\]
against
\[
A_{\mathrm{weighted}}
=
\sum_i w_i A_i,
\]
where the weights may depend on resolution, layer, head, or an uncertainty measure.

This experiment will be carried out after the simpler resolution and timestep
ablations, since their results can guide which attention maps should receive greater
importance.

\textbf{Comparison with existing diffusion-based segmentation methods.}
In addition to the ablations and method-specific extensions described above, we
will compare our approach against existing diffusion-based segmentation methods
using the released OVAM benchmark data, including the generated images, prompts,
seeds, and ground-truth masks. For methods already evaluated in the OVAM paper,
namely DAAM, Attn2Mask, Grounded Diffusion, and DatasetDM, we will report the
results published by OVAM. For methods not included in the OVAM comparison, we
will reproduce them on the same benchmark where feasible. The main methods
considered for reproduction are DiffSegmenter, DiffSeg, SeeDiff, FA-Seg,
DiffuMask, and FreeSeg-Diff. This experiment will provide a broader comparison
between our proposed strategies and other Stable-Diffusion-based segmentation
approaches under a common evaluation setting.

% =========================================================
\subsection*{Experiment summary}
% =========================================================

The planned experiments are summarized below. The first group mainly reuses stored
cross-attention and self-attention information and can therefore be evaluated through
offline ablations. The second group requires additional method-specific components or
optimization.

\begin{table}[h]
\centering
\small
\begin{tabular}{p{3.2cm}p{3.0cm}p{3.6cm}p{2.1cm}}
\toprule
\textbf{Strategy} &
\textbf{Main attention / refinement} &
\textbf{Main ablation} &
\textbf{Result} \\
\midrule

\multicolumn{4}{l}{\textbf{Attention-map experiments / easier ablations}} \\
\midrule

OVAM baseline
& Cross-attention
& Normalization, threshold, background handling
& \\[1.2em]

OVAM + self-attention refinement
& Cross-attention + self-attention
& SA resolution, refinement strength, single vs.\ multi-resolution
& \\[1.2em]

Resolution-specific cross-attention
& Cross-attention
& $16\times16$, $32\times32$, $64\times64$, and combinations
& \\[1.2em]

Denoising timestep selection
& Cross-attention
& Early, middle, late, individual, and all timesteps
& \\[1.2em]

Class-/scale-adaptive attention
& Class-dependent attention selection
& Global configuration vs.\ class-/scale-dependent configuration
& \\[1.2em]

\midrule
\multicolumn{4}{l}{\textbf{Method-specific / harder extensions}} \\
\midrule

Above ablations + token optimization
& Optimized class embeddings
& Standard vs.\ optimized tokens across the previous attention configurations
& \\[1.2em]

Above ablations + SAM refinement
& Diffusion attention + SAM
& Point vs.\ multi-point vs.\ box vs.\ coarse-mask prompting
& \\[1.2em]

SeeDiff
& CA seeds + multi-scale SA expansion
& Seed threshold, SA resolutions, single vs.\ progressive expansion, background refinement
& \\[1.2em]

Hierarchical / weighted attention fusion
& Multi-resolution / multi-layer CA/SA
& Uniform vs.\ weighted vs.\ uncertainty-/entropy-based fusion
& \\[1.2em]

Comparison with existing methods
& OVAM benchmark + reproduced SD-based methods
& Report from OVAM: DAAM, Attn2Mask, Grounded Diffusion, DatasetDM; reproduce: DiffSegmenter, DiffSeg, SeeDiff, FA-Seg, DiffuMask, FreeSeg-Diff
& \\[1.2em]

\bottomrule
\end{tabular}
\caption{Planned segmentation experiments, ordered from attention-map ablations
that can mainly reuse stored attention information to more method-specific extensions.}
\end{table}

The final evaluation will report overall mIoU together with per-class IoU.
Particular attention will be given to the difference between large spatially coherent
classes such as road, sky, vegetation, and building, and small or thin classes such as
person, pole, traffic light, and traffic sign.

\textbf{Final note.}
The proposed order is tentative and may change according to the results of the
earlier experiments. In particular, the resolution, timestep, and self-attention
ablations are intended to guide the design of the later methods rather than treating
all possible configurations equally.

\end{document}